Expected Shortfall (ES), the market-risk measure of the Basel framework, is not elicitable on its own, but it is jointly elicitable with Value-at-Risk (VaR). We use this fact to train deep sequence models directly on a strictly consistent scoring function for (VaR, ES), the FZ0 loss. We prove that pooling across assets on volatility-standardised returns is exactly equivalent to training on raw returns. We also design an output layer that rules out crossing risk forecasts by construction. In a twenty-year out-of-sample study of twelve global equity indices at the 1\%, 2.5\% and 5\% levels, we compare an LSTM, a Transformer and a GARCH--LSTM hybrid with nine benchmarks, including GARCH-EVT, filtered historical simulation and score-driven FZ-GAS models. A pooled LSTM that sees only past returns is statistically indistinguishable from the best econometric models and is the only single model in every model confidence set. The deep models are, moreover, complementary to the econometric ones. Adding them to an equal-weight combination of the econometric models lowers the FZ0 loss at every risk level (Diebold--Mariano $t$ between $-4.7$ and $-5.3$; significant in 10--11 of 12 markets, never significantly worse). The resulting combination is the best model overall and requires about 2\% less capital than GJR-GARCH with the same Basel traffic-light record. The deep models' weaknesses are calibration and sudden volatility spikes such as March 2020.
